\documentclass[11pt]{article}

\usepackage[a4paper,margin=1in]{geometry}
\usepackage{xurl}
\usepackage{hyperref}
\usepackage{parskip}
\emergencystretch=3em

\title{Nakafa Lab Note: Thalassemia Carrier Status and Malaria}
\author{Nakafa Lab}
\date{2026-04-26}

\begin{document}

\maketitle

\section*{Working conclusion}

Some thalassemia carrier states can provide partial protection against severe
malaria. This does not mean a carrier is immune to malaria, and it does not make
severe thalassemia beneficial for an affected patient.

\section*{Evidence}

The 2025 Thalassaemia International Federation guideline describes
beta-thalassemia as common in regions historically affected by malaria and links
high gene frequency to selective pressure from \textit{Plasmodium falciparum}.
It also separates usually asymptomatic heterozygous carriers from severe disease
states that arise when high-risk variants combine.

Alpha-thalassemia malaria literature supports the broader evolutionary pattern:
red-cell variants can persist when carriers gain a survival advantage under
malaria pressure.

\section*{Interpretation}

This finding should not be used to label a specific patient's illness as
punishment. A more careful framing is that genetic variants can arise and persist
across generations, sometimes because carriers had historical survival
advantages. Severe disease remains a medical burden that needs better treatment,
screening, and access.

\section*{Research implication}

The malaria connection keeps the program focused on red-cell biology:
globin-chain balance, fetal hemoglobin induction, oxidative stress, immune
clearance, and affordable ways to reduce transfusion burden.

\section*{Sources}

\begin{itemize}
  \item TIF 2025 TDT guideline, genetic basis chapter:
  \url{https://www.ncbi.nlm.nih.gov/books/NBK614253/}
  \item Effect of alpha-thalassemia on malaria and other diseases in children:
  \url{https://journals.plos.org/plosmedicine/article?id=10.1371/journal.pmed.0030158}
\end{itemize}

\end{document}
